\documentclass[11pt]{letter}
\usepackage[margin=0.75in]{geometry}
\usepackage{parskip}
\usepackage{hyperref}

\signature{Ning Zhou\\
Public R\&D Center, SuperMap Software Co., Ltd.\\
Beijing, China\\
\href{mailto:zhouning1@supermap.com}{zhouning1@supermap.com}}

\address{}

\begin{document}

\begin{letter}{The Editor \\ \textit{Scientific Reports}}

\opening{Dear Editor,}

We are pleased to submit our manuscript, ``Reproducible model-based planning for county-scale farmland consolidation in fragmented mountain landscapes,'' for consideration as an Article in \textit{Scientific Reports}.

The manuscript evaluates whether a learned transition model with model-predictive control can generate auditable cadastral scenarios faster than model-free reinforcement learning and classical operations-research baselines. It also exposes the constraints that determine whether a candidate can be deployed.

We report evidence from two real Chinese counties, Bishan District (Chongqing, $52{,}515$ parcels) and Neijiang Dongxing District (Sichuan, $76{,}376$ parcels). The evaluation also includes a seven-preset synthetic benchmark and a public-data Buchanan County mine-restoration boundary check. On the matched Bishan cadastre, contrastive MPC achieves a $1.6\times$ larger slope reduction than in-house reinforcement-learning baselines. Independent GIS recomputation verifies that the canonical Bishan and Neijiang deployment deltas are properties of the optimised cadastral geometry.

The transfer framing is deliberately technical rather than promotional. The unconstrained Bishan candidate reduces farmland slope by $-2.001\%$ and shifts $687$\,ha out of steep bands. It also loses $505$\,ha of cultivated area and is rejected as a final policy-compliant plan. A no-net-loss cultivated-area floor retains a $-0.675\%$ slope change, higher contiguity, and $29.8$\,ha more qualifying large-patch area. Neijiang shows the same repair pattern, with $330.4$\,ha less farmland in $\geq 15^\circ$ slope bands under the constrained canonical audit. The revised Supplementary Information now includes seven-profile execution frontiers for both Bishan and Neijiang, while still avoiding a regional-generalisation claim. We therefore present an auditable scenario generator with explicit deployment boundaries, not a substitute for local plan approval.

The work is suitable for \textit{Scientific Reports} because its main contribution is technical validity, reproducibility, and bounded empirical evaluation. The evidence spans real, synthetic, and public-data settings. The release includes trained ensembles, an ArcGIS Pro Python Toolbox, a command-line workflow, a smoke test, and open reproduction tracks. These tracks cover the synthetic benchmark, the Buchanan boundary check, and the training/planning diagnostics. Derived data products support verification of the real-county results without redistributing raw cadastral geometries. Code and reproduction instructions are available at \href{https://github.com/zhouning/arcgis-farmland-mpc}{github.com/zhouning/arcgis-farmland-mpc}.

The manuscript is not under consideration elsewhere. The authors declare no competing interests, do not request exclusion of any referees, and have had no prior discussions with \textit{Scientific Reports} about this submission.

\closing{Sincerely,}

\end{letter}
\end{document}
